\DefineFigure{pic0201}{
\begin{figure}[H]
  \vspace{-0.3cm}
  \centering
  \begin{notescircuit}[scale=0.9]
    \ctikzset{sources/symbol/sign rotation=straight}
    \draw (0,2.0) node[ocirc] {} -- (2.3,2.0) to[american controlled current source,l_=$M(D)I_{\mathrm{out}}$] (2.3,0) -- (0,0) node[ocirc] {};
    \draw[-{Latex[length=1.7mm]}] (0.4,2.0) -- (1.3,2.0) node[midway,above] {$I_{\mathrm{in}}$};
    \node at (-0.35,2.0) {$+$};
    \node at (-0.35,0) {$-$};
    \node at (-0.35,1.0) {$V_{\mathrm{in}}$};
    \draw (3.67,2.0) to[american controlled voltage source,l=$M(D)V_{\mathrm{in}}$] (3.67,0);
    \draw (3.67,2.0) -- (6.17,2.0) node[ocirc] {};
    \draw (3.67,0) -- (6.17,0) node[ocirc] {};
    \draw[-{Latex[length=1.7mm]}] (4.47,2.0) -- (5.47,2.0) node[midway,above] {$I_{\mathrm{out}}$};
    \node at (1.1,-0.55) {功率输入};
    \node at (4.97,-0.55) {功率输出};
    \node at (6.52,2.0) {$+$};
    \node at (6.52,0) {$-$};
    \node at (6.52,1.0) {$V_{\mathrm{out}}$};
  \end{notescircuit}
  \vspace{-0.15cm}
  \caption{理想变换器的受控源稳态模型}
  \label{pic0201}
  \vspace{-0.3cm}
\end{figure}
}

\DefineFigure{pic0202}{
\begin{figure}[H]
  \vspace{-0.3cm}
  \centering
  \begin{subfigure}[t]{0.48\textwidth}
    \centering
    \begin{notescircuit}[scale=0.74]
      \node[transformer] (tr) at (3.0,1.3) {};
      \node[circ] at (tr.outer dot A1) {};
      \node[circ] at (tr.outer dot B1) {};
      \node[above] at (3.0,3.05) {$1:M(D)$};
      \coordinate (dc-mark) at ($(tr-L1.a)!{0.5/\ctikzvalof{bipoles/americaninductor/coils}}!(tr-L1.b)$);
      \draw (tr.outer dot A1 |- dc-mark) -- (tr.outer dot B1 |- dc-mark);
      \draw (0,2.6) to[european voltage source,l=$V_{\mathrm{in}}$,name=source-primary] (0,0);
      \node[above left=1pt] at (source-primary.b) {$+$};
      \node[below left=1pt] at (source-primary.a) {$-$};
      \draw (0,2.6) to[european resistor,l=$R_1$] (1.3,2.6) -- (2.0,2.6) -| (tr.A1);
      \draw (tr.A2) |- (0,0);
      \draw (tr.B1) |- (5.3,2.6);
      \draw (5.3,2.6) to[european resistor,l_=$R$] (5.3,0) -| (tr.B2);
      \node at (5.85,2.6) {$+$};
      \node at (5.85,0) {$-$};
      \node at (5.85,1.3) {$V_{\mathrm{out}}$};
      \draw[-{Latex[length=1.7mm]}] (1.45,2.6) -- (2.15,2.6) node[midway,above] {$I_{\mathrm{in}}$};
      \draw[-{Latex[length=1.7mm]}] (3.0,-0.35) -- ([yshift=-0.05cm]tr.south) node[midway,left] {$D$};
    \end{notescircuit}
    \caption{原边模型}
    \label{pic0202a}
  \end{subfigure}
  \hfill
  \begin{subfigure}[t]{0.48\textwidth}
    \centering
    \begin{notescircuit}[scale=0.83]
      \draw (0,2.6) to[european voltage source,l=$M(D)V_{\mathrm{in}}$,name=source-secondary] (0,0);
      \node[above left=1pt] at (source-secondary.b) {$+$};
      \node[below left=1pt] at (source-secondary.a) {$-$};
      \draw (0,2.6) to[european resistor,l=$M^2(D)R_1$] (3.6,2.6)
      to[european resistor,l_=$R$] (3.6,0) -- (0,0);
      \node at (4.15,2.6) {$+$};
      \node at (4.15,0) {$-$};
      \node at (4.15,1.3) {$V_{\mathrm{out}}$};
    \end{notescircuit}
    \caption{折算到副边}
    \label{pic0202b}
  \end{subfigure}
  \vspace{-0.15cm}
  \caption{直流“变压器”模型及阻抗折算}
  \label{pic0202}
  \vspace{-0.3cm}
\end{figure}
}

\DefineFigure{pic0203}{
\begin{figure}[H]
  \vspace{-0.5cm}
  \centering
  \begin{subfigure}[t]{0.48\textwidth}
    \centering
    \begin{notescircuit}[scale=0.9]
      \ctikzset{
        tripoles/nigfete/base width=0.32,
        tripoles/nigfete/gate width=0.40
      }
      \coordinate (gnd-rl-a) at (0,0);
      \node[nigfete, nogate, anchor=center] (q-rl-a) at (2.98,0.95) {};
      \draw (gnd-rl-a)
      to[V, l_=$V_{\mathrm{in}}$, name=vin-rl-a] (0,2.0)
      -- (0.57,2.0)
      to[L, l=$L$, name=ind-rl-a] (1.68,2.0) coordinate (ind-end-rl-a)
      to[european resistor, l=$R_{\mathrm{L}}$] (2.98,2.0) coordinate (sw-rl-a)
      to[D-, l=$D$] (4.28,2.0) coordinate (out-rl-a)
      -- ++(1.05,0) coordinate (load-rl-a);
      \draw[-{Latex[length=1.5mm]}] (0.12,2.0) -- (0.52,2.0) node[midway,above] {$i_{\mathrm{L}}$};
      \draw (sw-rl-a) -- (q-rl-a.D);
      \draw (q-rl-a.S) -- (q-rl-a.S |- gnd-rl-a);
      \draw (q-rl-a.inner down -| q-rl-a.nogate) -- ++(-0.22,0);
      \node[left=11pt] at (q-rl-a.center) {$Q$};
      \node[left=1pt] at (0,1.35) {$+$};
      \node[left=1pt] at (0,0.65) {$-$};
      \node[below=4pt] at (ind-rl-a.center)
      {$+\;v_{\mathrm{L}}\;-$};
      \draw (out-rl-a) to[C, l_=$C$] (out-rl-a |- gnd-rl-a);
      \draw (load-rl-a) to[european resistor, l_=$R$] (load-rl-a |- gnd-rl-a);
      \draw (load-rl-a |- gnd-rl-a) -- (gnd-rl-a);
      \node[right] at (load-rl-a) {$+$};
      \node[right] at (load-rl-a |- gnd-rl-a) {$-$};
      \path (load-rl-a) -- node[right] {$v_{\mathrm{out}}$} (load-rl-a |- gnd-rl-a);
    \end{notescircuit}
    \caption{原电路}
    \label{pic0203a}
  \end{subfigure}
  \hfill
  \begin{subfigure}[t]{0.48\textwidth}
    \centering
    \begin{notescircuit}[scale=0.9]
      \coordinate (gnd-rl-b) at (0,0);
      \coordinate (ind-end-rl-b) at (1.61,2.0);
      \coordinate (pole-rl-b) at (2.91,2.0);
      \coordinate (throw-one-rl-b) at (3.31,1.25);
      \coordinate (throw-two-rl-b) at (3.66,2.0);
      \draw (gnd-rl-b)
      to[V, l_=$V_{\mathrm{in}}$, name=vin-rl-b] (0,2.0)
      -- (0.54,2.0)
      to[L, l=$L$, name=ind-rl-b] (ind-end-rl-b)
      to[european resistor, l=$R_{\mathrm{L}}$] (pole-rl-b);
      \draw[-{Latex[length=1.5mm]}] (0.12,2.0) -- (0.52,2.0) node[midway,above] {$i_{\mathrm{L}}$};
      \node[ocirc] at (pole-rl-b) {};
      \node[ocirc] at (throw-one-rl-b) {};
      \node[ocirc] at (throw-two-rl-b) {};
      \draw (pole-rl-b) -- (3.54,1.60);
      \draw (throw-one-rl-b) -- (throw-one-rl-b |- gnd-rl-b) -- (gnd-rl-b);
      \draw (throw-two-rl-b) -- ++(0.65,0) coordinate (out-rl-b)
      -- ++(1.05,0) coordinate (load-rl-b);
      \node[left] at (throw-one-rl-b) {$1$};
      \node[above] at (throw-two-rl-b) {$2$};
      \node[left=1pt] at (0,1.35) {$+$};
      \node[left=1pt] at (0,0.65) {$-$};
      \node[below=4pt] at (ind-rl-b.center)
      {$+\;v_{\mathrm{L}}\;-$};
      \draw (out-rl-b) to[C, l_=$C$] (out-rl-b |- gnd-rl-b);
      \draw (load-rl-b) to[european resistor, l_=$R$] (load-rl-b |- gnd-rl-b);
      \draw (load-rl-b |- gnd-rl-b) -- (gnd-rl-b);
      \node[right] at (load-rl-b) {$+$};
      \node[right] at (load-rl-b |- gnd-rl-b) {$-$};
      \path (load-rl-b) -- node[right] {$v_{\mathrm{out}}$} (load-rl-b |- gnd-rl-b);
    \end{notescircuit}
    \caption{等效电路}
    \label{pic0203b}
  \end{subfigure}
  \medskip
  \begin{subfigure}[t]{0.48\textwidth}
    \centering
    \begin{notescircuit}[scale=0.9]
      \coordinate (gnd-rl-c) at (0,0);
      \draw (gnd-rl-c)
      to[V, l_=$V_{\mathrm{in}}$, name=vin-rl-c] (0,1.8)
      -- (0.56,1.8)
      to[L, l=$L$, name=ind-rl-c] (1.64,1.8)
      to[european resistor, l=$R_{\mathrm{L}}$] (2.94,1.8)
      -- (2.94,0);
      \draw[-{Latex[length=1.5mm]}] (0.12,1.8) -- (0.52,1.8) node[midway,above] {$i_{\mathrm{L}}$};
      \draw (3.84,1.8) coordinate (out-rl-c)
      -- ++(1.05,0) coordinate (load-rl-c);
      \draw (out-rl-c) to[C, l_=$C$] (out-rl-c |- gnd-rl-c);
      \draw (load-rl-c) to[european resistor, l_=$R$] (load-rl-c |- gnd-rl-c);
      \draw (load-rl-c |- gnd-rl-c) -- (gnd-rl-c);
      \node[left=1pt] at (0,1.25) {$+$};
      \node[left=1pt] at (0,0.55) {$-$};
      \node[below=4pt] at (ind-rl-c.center)
      {$+\;v_{\mathrm{L}}\;-$};
      \node[right] at (load-rl-c) {$+$};
      \node[right] at (load-rl-c |- gnd-rl-c) {$-$};
      \path (load-rl-c) -- node[right] {$v_{\mathrm{out}}$} (load-rl-c |- gnd-rl-c);
    \end{notescircuit}
    \caption{状态 1}
    \label{pic0203c}
  \end{subfigure}
  \hfill
  \begin{subfigure}[t]{0.48\textwidth}
    \centering
    \begin{notescircuit}[scale=0.9]
      \coordinate (gnd-rl-d) at (0,0);
      \draw (gnd-rl-d)
      to[V, l_=$V_{\mathrm{in}}$, name=vin-rl-d] (0,1.8)
      -- (0.64,1.8)
      to[L, l=$L$, name=ind-rl-d] (1.91,1.8)
      to[european resistor, l=$R_{\mathrm{L}}$] (3.21,1.8) coordinate (out-rl-d)
      -- ++(1.05,0) coordinate (load-rl-d);
      \draw[-{Latex[length=1.5mm]}] (0.12,1.8) -- (0.52,1.8) node[midway,above] {$i_{\mathrm{L}}$};
      \draw (out-rl-d) to[C, l_=$C$] (out-rl-d |- gnd-rl-d);
      \draw (load-rl-d) to[european resistor, l_=$R$] (load-rl-d |- gnd-rl-d);
      \draw (load-rl-d |- gnd-rl-d) -- (gnd-rl-d);
      \node[left=1pt] at (0,1.25) {$+$};
      \node[left=1pt] at (0,0.55) {$-$};
      \node[below=4pt] at (ind-rl-d.center)
      {$+\;v_{\mathrm{L}}\;-$};
      \node[right] at (load-rl-d) {$+$};
      \node[right] at (load-rl-d |- gnd-rl-d) {$-$};
      \path (load-rl-d) -- node[right] {$v_{\mathrm{out}}$} (load-rl-d |- gnd-rl-d);
    \end{notescircuit}
    \caption{状态 2}
    \label{pic0203d}
  \end{subfigure}
  \vspace{-0.5cm}
  \caption{含电感铜损的 Boost 变换器及其开关状态等效电路}
  \vspace{-0.5cm}
  \label{pic0203}
\end{figure}
}

\DefineFigure{pic0204}{
\begin{figure}[H]
  \centering
  \begin{tikzpicture}[
    axis/.style={-{Latex[length=2.2mm]}, line width=0.6pt},
    guide/.style={densely dashed, draw=black, line width=0.4pt},
    waveform/.style={line width=0.9pt},
    every node/.style={font=\small}
    ]
    \def\xzero{1.5}
    \def\xduty{5.10}
    \def\xperiod{10.5}
    \def\xend{11.8}
    \def\ytiming{5.85}
    % 两个开关状态的时间分段；竖直参考线贯穿上下两幅波形
    \draw[guide] (\xduty,0.35) -- (\xduty,\ytiming);
    \draw[guide] (\xperiod,0.35) -- (\xperiod,\ytiming);
    \draw[{Latex[length=1.6mm]}-{Latex[length=1.6mm]}, line width=0.5pt]
    (\xzero,\ytiming) -- node[above] {$DT_{\mathrm{s}}$} (\xduty,\ytiming);
    \draw[{Latex[length=1.6mm]}-{Latex[length=1.6mm]}, line width=0.5pt]
    (\xduty,\ytiming) -- node[above] {$D'T_{\mathrm{s}}$} (\xperiod,\ytiming);
    % Boost 电感电流纹波
    \draw[axis] (\xzero,3.55) -- (\xend,3.55) node[right] {$t$};
    \draw[axis] (\xzero,3.55) -- (\xzero,6.25)
    node[left] {$i_{\mathrm{L}}(t)$};
    \node[left] at (\xzero,3.55) {$0$};
    \draw[guide] (\xzero,4.85) -- (\xperiod,4.85);
    \node[left] at (\xzero,4.85) {$I_{\mathrm{L}}$};
    \draw[waveform]
    (\xzero,4.20) -- (\xduty,5.50) -- (\xperiod,4.20) -- (\xend,{4.20+(5.50-4.20)*(\xend-\xperiod)/(\xduty-\xzero)});
    \node at ({(\xzero+\xduty)/2},4.00)
    {$i'_{\mathrm{L}}(t)\approx\dfrac{V_{\mathrm{in}}-I_{\mathrm{L}}R_{\mathrm{L}}}{L}$};
    \node at ({(\xduty+\xperiod)/2},4.00)
    {$i'_{\mathrm{L}}(t)\approx\dfrac{V_{\mathrm{in}}-I_{\mathrm{L}}R_{\mathrm{L}}-V_{\mathrm{out}}}{L}$};
    \draw[guide] (\xduty,5.50) -- (11.05,5.50);
    \draw[guide] (\xperiod,4.20) -- (11.05,4.20);
    \draw[{Latex[length=1.7mm]}-{Latex[length=1.7mm]}, line width=0.6pt]
    (11.05,4.20) -- node[right] {$\Delta i_{\mathrm{L,pp}}$} (11.05,5.50);
    % Boost 电容电压纹波
    \draw[axis] (\xzero,0.35) -- (\xend,0.35) node[right] {$t$};
    \draw[axis] (\xzero,0.35) -- (\xzero,2.85)
    node[left] {$v_{\mathrm{C}}(t)$};
    \node[left] at (\xzero,0.35) {$0$};
    \draw[guide] (\xzero,1.55) -- (\xperiod,1.55);
    \node[left] at (\xzero,1.55) {$V_{\mathrm{out}}$};
    \draw[waveform]
    (\xzero,1.95) -- (\xduty,1.15) -- (\xperiod,1.95) -- (\xend,{1.95-(1.95-1.15)*(\xend-\xperiod)/(\xduty-\xzero)});
    \node at ({(\xzero+\xduty)/2},2.35)
    {$v'_{\mathrm{C}}(t)\approx-\dfrac{V_{\mathrm{out}}}{RC}$};
    \node at ({(\xduty+\xperiod)/2},2.35)
    {$v'_{\mathrm{C}}(t)\approx\dfrac{I_{\mathrm{L}}}{C}-\dfrac{V_{\mathrm{out}}}{RC}$};
    \draw[guide] (\xduty,1.15) -- (11.05,1.15);
    \draw[guide] (\xperiod,1.95) -- (11.05,1.95);
    \draw[{Latex[length=1.7mm]}-{Latex[length=1.7mm]}, line width=0.6pt]
    (11.05,1.15) -- node[right] {$\Delta v_{\mathrm{C,pp}}$} (11.05,1.95);
    \node[below] at (\xperiod,0.35) {$T_{\mathrm{s}}$};
  \end{tikzpicture}
  \vspace{-0.2cm}
  \caption{含电感铜损的 Boost 变换器电感电流与电容电压纹波}
  \vspace{-0.5cm}
  \label{pic0204}
\end{figure}
}

\DefineFigure{pic0205}{
\begin{figure}[H]
  \centering
  \begin{subfigure}[t]{0.48\textwidth}
    \centering
    \begin{notescircuit}[scale=0.7]
      \ctikzset{sources/symbol/sign rotation=straight}
      \draw (0,2.6) to[european voltage source,l=$V_{\mathrm{in}}$,name=boost-source-model] (0,0);
      \node[above left=1pt] at (boost-source-model.b) {$+$};
      \node[below left=1pt] at (boost-source-model.a) {$-$};
      \draw (0,2.6) -- (0.9,2.6)
      to[european resistor,l=$R_{\mathrm{L}}$] (2.2,2.6) -- (3.7,2.6);
      \draw (3.7,2.6) to[american controlled voltage source,l_=$D'V_{\mathrm{out}}$] (3.7,0);
      \draw (3.7,0) -- (0,0);
      \draw[-{Latex[length=1.7mm]}] (2.35,2.6) -- (3.05,2.6) node[midway,above] {$I_{\mathrm{L}}$};
      \draw (5.1,0) to[american controlled current source,l_=$D'I_{\mathrm{L}}$] (5.1,2.6)
      -- (7.6,2.6) to[european resistor,l_=$R$] (7.6,0) -- (5.1,0);
      \draw[-{Latex[length=1.7mm]}] (5.6,2.6) -- (6.5,2.6) node[midway,above] {$I_{\mathrm{out}}$};
      \node at (8.15,2.6) {$+$};
      \node at (8.15,0) {$-$};
      \node at (8.15,1.3) {$V_{\mathrm{out}}$};
    \end{notescircuit}
    \caption{受控源模型}
    \label{pic0205a}
  \end{subfigure}
  \hfill
  \begin{subfigure}[t]{0.48\textwidth}
    \centering
    \begin{notescircuit}[scale=0.74]
      \node[transformer] (tr) at (3.0,1.3) {};
      \node[circ] at (tr.outer dot A1) {};
      \node[circ] at (tr.outer dot B1) {};
      \node[above] at (3.0,3.05) {$D':1$};
      \coordinate (boost-dc-mark) at ($(tr-L1.a)!{0.5/\ctikzvalof{bipoles/americaninductor/coils}}!(tr-L1.b)$);
      \draw (tr.outer dot A1 |- boost-dc-mark) -- (tr.outer dot B1 |- boost-dc-mark);
      \draw (0,2.6) to[european voltage source,l=$V_{\mathrm{in}}$,name=boost-source-transformer] (0,0);
      \node[above left=1pt] at (boost-source-transformer.b) {$+$};
      \node[below left=1pt] at (boost-source-transformer.a) {$-$};
      \draw (0,2.6) to[european resistor,l=$R_{\mathrm{L}}$] (1.3,2.6) -- (2.0,2.6) -| (tr.A1);
      \draw (tr.A2) |- (0,0);
      \draw (tr.B1) |- (5.3,2.6);
      \draw (5.3,2.6) to[european resistor,l_=$R$] (5.3,0) -| (tr.B2);
      \node at (5.85,2.6) {$+$};
      \node at (5.85,0) {$-$};
      \node at (5.85,1.3) {$V_{\mathrm{out}}$};
      \draw[-{Latex[length=1.7mm]}] (1.45,2.6) -- (2.15,2.6) node[midway,above] {$I_{\mathrm{L}}$};
      \draw[-{Latex[length=1.7mm]}] (4.0,2.6) -- (4.7,2.6) node[midway,above] {$I_{\mathrm{out}}$};
      \draw[-{Latex[length=1.7mm]}] (3.0,-0.35) -- ([yshift=-0.05cm]tr.south) node[midway,left] {$D$};
    \end{notescircuit}
    \caption{直流“变压器”模型}
    \label{pic0205b}
  \end{subfigure}
  \vspace{-0.15cm}
  \caption{Boost 变换器的直流稳态等效模型}
  \label{pic0205}
\end{figure}
}

\DefineFigure{pic0206}{
\begin{figure}[H]
  \centering
  \begin{notescircuit}[scale=0.95]
    \draw (0,2.6) to[european voltage source,l=$V_{\mathrm{in}}/D'$,name=boost-source-reflected] (0,0);
    \node[above left=1pt] at (boost-source-reflected.b) {$+$};
    \node[below left=1pt] at (boost-source-reflected.a) {$-$};
    \draw (0,2.6) -- (0.4,2.6)
    to[european resistor,l=$R_{\mathrm{L}}/(D')^2$] (1.7,2.6) -- (4.0,2.6)
    to[european resistor,l_=$R$] (4.0,0) -- (0,0);
    \draw[-{Latex[length=1.7mm]}] (2.1,2.6) -- (3.2,2.6) node[midway,above] {$D'I_{\mathrm{L}}$};
    \node at (4.55,2.6) {$+$};
    \node at (4.55,0) {$-$};
    \node at (4.55,1.3) {$V_{\mathrm{out}}$};
  \end{notescircuit}
  \vspace{-0.15cm}
  \caption{Boost 稳态模型折算到输出侧}
  \label{pic0206}
\end{figure}
}

\DefineFigure{pic0207}{
\begin{figure}[H]
  \vspace{-0.2cm}
  \centering
  \begin{subfigure}[t]{0.48\textwidth}
    \centering
    \begin{notescircuit}[scale=0.9]
      \ctikzset{
        tripoles/nigfete/base width=0.32,
        tripoles/nigfete/gate width=0.40
      }
      \coordinate (gnd-buck-rl-a) at (0,0);
      \node[nigfete, nogate, rotate=90, anchor=center] (q-buck-rl-a) at (1.05,2.0) {};
      \draw (gnd-buck-rl-a)
      to[V, l_=$V_{\mathrm{in}}$, name=vin-buck-rl-a] (0,2.0)
      -- (q-buck-rl-a.D);
      \draw (q-buck-rl-a.inner down |- q-buck-rl-a.nogate) -- ++(0,-0.22);
      \node[above] at (q-buck-rl-a.center) {$Q$};
      \draw (q-buck-rl-a.S) -- ++(0.35,0) coordinate (sw-buck-rl-a)
      -- ++(0.56,0)
      to[L, l=$L$, name=ind-buck-rl-a] ++(1.08,0)
      to[european resistor, l=$R_{\mathrm{L}}$] ++(1.30,0) coordinate (out-buck-rl-a)
      -- ++(1.05,0) coordinate (load-buck-rl-a);
      \draw[-{Latex[length=1.5mm]}] ($(sw-buck-rl-a)+(0.12,0)$) -- ($(sw-buck-rl-a)+(0.52,0)$)
      node[midway,above] {$i_{\mathrm{L}}$};
      \draw (sw-buck-rl-a |- gnd-buck-rl-a) to[D-, l_=$D$] (sw-buck-rl-a);
      \node[left=1pt] at (0,1.35) {$+$};
      \node[left=1pt] at (0,0.65) {$-$};
      \node[below=4pt] at (ind-buck-rl-a.center)
      {$+\;v_{\mathrm{L}}\;-$};
      \draw (out-buck-rl-a) to[C, l_=$C$] (out-buck-rl-a |- gnd-buck-rl-a);
      \draw (load-buck-rl-a) to[european resistor, l_=$R$] (load-buck-rl-a |- gnd-buck-rl-a);
      \draw (load-buck-rl-a |- gnd-buck-rl-a) -- (gnd-buck-rl-a);
      \node[right] at (load-buck-rl-a) {$+$};
      \node[right] at (load-buck-rl-a |- gnd-buck-rl-a) {$-$};
      \path (load-buck-rl-a) -- node[right] {$v_{\mathrm{out}}$} (load-buck-rl-a |- gnd-buck-rl-a);
    \end{notescircuit}\par
    \caption{原电路}
    \label{pic0207a}
  \end{subfigure}
  \hfill
  \begin{subfigure}[t]{0.48\textwidth}
    \centering
    \begin{notescircuit}[scale=0.9]
      \coordinate (gnd-buck-rl-b) at (0,0);
      \coordinate (throw-one-buck-rl-b) at (0.95,2.0);
      \coordinate (throw-two-buck-rl-b) at (1.33,1.25);
      \coordinate (pole-buck-rl-b) at (1.7,2.0);
      \draw (gnd-buck-rl-b)
      to[V, l_=$V_{\mathrm{in}}$, name=vin-buck-rl-b] (0,2.0)
      -- (throw-one-buck-rl-b);
      \draw (throw-two-buck-rl-b) -- (throw-two-buck-rl-b |- gnd-buck-rl-b) -- (gnd-buck-rl-b);
      \node[ocirc] at (throw-one-buck-rl-b) {};
      \node[ocirc] at (throw-two-buck-rl-b) {};
      \node[ocirc] at (pole-buck-rl-b) {};
      \draw (pole-buck-rl-b) -- (1.08,1.58);
      \node[above] at (throw-one-buck-rl-b) {$1$};
      \node[left] at (throw-two-buck-rl-b) {$2$};
      \draw (pole-buck-rl-b) -- ++(0.35,0) coordinate (sw-buck-rl-b)
      -- ++(0.56,0)
      to[L, l=$L$, name=ind-buck-rl-b] ++(1.08,0)
      to[european resistor, l=$R_{\mathrm{L}}$] ++(1.30,0) coordinate (out-buck-rl-b)
      -- ++(1.05,0) coordinate (load-buck-rl-b);
      \draw[-{Latex[length=1.5mm]}] ($(sw-buck-rl-b)+(0.12,0)$) -- ($(sw-buck-rl-b)+(0.52,0)$)
      node[midway,above] {$i_{\mathrm{L}}$};
      \node[left=1pt] at (0,1.35) {$+$};
      \node[left=1pt] at (0,0.65) {$-$};
      \node[below=4pt] at (ind-buck-rl-b.center)
      {$+\;v_{\mathrm{L}}\;-$};
      \draw (out-buck-rl-b) to[C, l_=$C$] (out-buck-rl-b |- gnd-buck-rl-b);
      \draw (load-buck-rl-b) to[european resistor, l_=$R$] (load-buck-rl-b |- gnd-buck-rl-b);
      \draw (load-buck-rl-b |- gnd-buck-rl-b) -- (gnd-buck-rl-b);
      \node[right] at (load-buck-rl-b) {$+$};
      \node[right] at (load-buck-rl-b |- gnd-buck-rl-b) {$-$};
      \path (load-buck-rl-b) -- node[right] {$v_{\mathrm{out}}$} (load-buck-rl-b |- gnd-buck-rl-b);
    \end{notescircuit}\par
    \caption{等效电路}
    \label{pic0207b}
  \end{subfigure}
  \medskip
  \begin{subfigure}[t]{0.48\textwidth}
    \centering
    \begin{notescircuit}[scale=0.9]
      \coordinate (gnd-buck-rl-c) at (0,0);
      \draw (gnd-buck-rl-c)
      to[V, l_=$V_{\mathrm{in}}$, name=vin-buck-rl-c] (0,1.8)
      -- (0.56,1.8)
      to[L, l=$L$, name=ind-buck-rl-c] (1.64,1.8)
      to[european resistor, l=$R_{\mathrm{L}}$] (2.94,1.8) coordinate (out-buck-rl-c)
      -- ++(1.05,0) coordinate (load-buck-rl-c);
      \draw[-{Latex[length=1.5mm]}] (0.12,1.8) -- (0.52,1.8) node[midway,above] {$i_{\mathrm{L}}$};
      \node[left=1pt] at (0,1.25) {$+$};
      \node[left=1pt] at (0,0.55) {$-$};
      \node[below=4pt] at (ind-buck-rl-c.center)
      {$+\;v_{\mathrm{L}}\;-$};
      \draw (out-buck-rl-c) to[C, l_=$C$] (out-buck-rl-c |- gnd-buck-rl-c);
      \draw (load-buck-rl-c) to[european resistor, l_=$R$] (load-buck-rl-c |- gnd-buck-rl-c);
      \draw (load-buck-rl-c |- gnd-buck-rl-c) -- (gnd-buck-rl-c);
      \node[right] at (load-buck-rl-c) {$+$};
      \node[right] at (load-buck-rl-c |- gnd-buck-rl-c) {$-$};
      \path (load-buck-rl-c) -- node[right] {$v_{\mathrm{out}}$} (load-buck-rl-c |- gnd-buck-rl-c);
    \end{notescircuit}
    \caption{状态 1}
    \label{pic0207c}
  \end{subfigure}
  \hfill
  \begin{subfigure}[t]{0.48\textwidth}
    \centering
    \begin{notescircuit}[scale=0.9]
      \coordinate (gnd-buck-rl-d) at (0,0);
      \draw (gnd-buck-rl-d) -- (0,1.8) -- (0.56,1.8)
      to[L, l=$L$, name=ind-buck-rl-d] (1.64,1.8)
      to[european resistor, l=$R_{\mathrm{L}}$] (2.94,1.8) coordinate (out-buck-rl-d)
      -- ++(1.05,0) coordinate (load-buck-rl-d);
      \draw[-{Latex[length=1.5mm]}] (0.12,1.8) -- (0.52,1.8) node[midway,above] {$i_{\mathrm{L}}$};
      \node[below=4pt] at (ind-buck-rl-d.center)
      {$+\;v_{\mathrm{L}}\;-$};
      \draw (out-buck-rl-d) to[C, l_=$C$] (out-buck-rl-d |- gnd-buck-rl-d);
      \draw (load-buck-rl-d) to[european resistor, l_=$R$] (load-buck-rl-d |- gnd-buck-rl-d);
      \draw (load-buck-rl-d |- gnd-buck-rl-d) -- (gnd-buck-rl-d);
      \node[right] at (load-buck-rl-d) {$+$};
      \node[right] at (load-buck-rl-d |- gnd-buck-rl-d) {$-$};
      \path (load-buck-rl-d) -- node[right] {$v_{\mathrm{out}}$} (load-buck-rl-d |- gnd-buck-rl-d);
    \end{notescircuit}
    \caption{状态 2}
    \label{pic0207d}
  \end{subfigure}
  \vspace{-0.5cm}
  \caption{含电感铜损的 Buck 变换器及其开关状态等效电路}
  \vspace{-0.7cm}
  \label{pic0207}
\end{figure}
}

\DefineFigure{pic0208}{
\begin{figure}[H]
  \centering
  \begin{subfigure}[t]{0.48\textwidth}
    \centering
    \begin{notescircuit}[scale=0.7]
      \ctikzset{sources/symbol/sign rotation=straight}
      \draw (0,2.6) to[european voltage source,l=$V_{\mathrm{in}}$,name=buck-source-model] (0,0);
      \node[above left=1pt] at (buck-source-model.b) {$+$};
      \node[below left=1pt] at (buck-source-model.a) {$-$};
      \draw (0,2.6) -- (2.8,2.6)
      to[american controlled current source,l_=$DI_{\mathrm{L}}$] (2.8,0) -- (0,0);
      \draw[-{Latex[length=1.7mm]}] (1.45,2.6) -- (2.15,2.6) node[midway,above] {$I_{\mathrm{in}}$};
      \draw (4.2,2.6) to[american controlled voltage source,l=$DV_{\mathrm{in}}$] (4.2,0);
      \draw (4.2,2.6) -- (5.1,2.6)
      to[european resistor,l=$R_{\mathrm{L}}$] (6.4,2.6) -- (7.6,2.6)
      to[european resistor,l_=$R$] (7.6,0) -- (4.2,0);
      \draw[-{Latex[length=1.7mm]}] (6.55,2.6) -- (7.25,2.6) node[midway,above] {$I_{\mathrm{L}}$};
      \node at (8.15,2.6) {$+$};
      \node at (8.15,0) {$-$};
      \node at (8.15,1.3) {$V_{\mathrm{out}}$};
    \end{notescircuit}
    \caption{受控源模型}
    \label{pic0208a}
  \end{subfigure}
  \hfill
  \begin{subfigure}[t]{0.48\textwidth}
    \centering
    \begin{notescircuit}[scale=0.74]
      \node[transformer] (tr) at (3.0,1.3) {};
      \node[circ] at (tr.outer dot A1) {};
      \node[circ] at (tr.outer dot B1) {};
      \node[above] at (3.0,3.05) {$1:D$};
      \coordinate (buck-dc-mark) at ($(tr-L1.a)!{0.5/\ctikzvalof{bipoles/americaninductor/coils}}!(tr-L1.b)$);
      \draw (tr.outer dot A1 |- buck-dc-mark) -- (tr.outer dot B1 |- buck-dc-mark);
      \draw (0,2.6) to[european voltage source,l=$V_{\mathrm{in}}$,name=buck-source-transformer] (0,0);
      \node[above left=1pt] at (buck-source-transformer.b) {$+$};
      \node[below left=1pt] at (buck-source-transformer.a) {$-$};
      \draw (0,2.6) -- (2.0,2.6) -| (tr.A1);
      \draw (tr.A2) |- (0,0);
      \draw (tr.B1) |- (4.0,2.6) to[european resistor,l=$R_{\mathrm{L}}$] (5.3,2.6);
      \draw (5.3,2.6) to[european resistor,l_=$R$] (5.3,0) -| (tr.B2);
      \node at (5.85,2.6) {$+$};
      \node at (5.85,0) {$-$};
      \node at (5.85,1.3) {$V_{\mathrm{out}}$};
      \draw[-{Latex[length=1.7mm]}] (1.45,2.6) -- (2.15,2.6) node[midway,above] {$I_{\mathrm{in}}$};
      \draw[-{Latex[length=1.7mm]}] (3.8,2.6) -- (4.2,2.6) node[midway,above] {$I_{\mathrm{L}}$};
      \draw[-{Latex[length=1.7mm]}] (3.0,-0.35) -- ([yshift=-0.05cm]tr.south) node[midway,left] {$D$};
    \end{notescircuit}
    \caption{直流“变压器”模型}
    \label{pic0208b}
  \end{subfigure}
  \vspace{-0.15cm}
  \caption{含电感铜损的 Buck 变换器直流稳态等效模型}
  \vspace{-0.4cm}
  \label{pic0208}
\end{figure}
}

\DefineFigure{pic0209}{
\begin{figure}[H]
  \centering
  \begin{notescircuit}[scale=0.9]
    \ctikzset{
      tripoles/nigfete/base width=0.32,
      tripoles/nigfete/gate width=0.40
    }
    \draw (0,0) to[V,l_=$V_{\mathrm{in}}$] (0,2.7)
    to[L,l=$L$,i>^=$i_{\mathrm{L}}$] (3.3,2.7)
    to[D-,l=$D$,name=diode] (6.0,2.7) -- (7.5,2.7);
    \node[nigfete, nogate, anchor=center] (q) at (3.3,1.25) {};
    \draw (3.3,2.7) -- (q.D);
    \draw (q.S) -- (q.S |- 0,0);
    \draw (q.inner down -| q.nogate) -- ++(-0.22,0);
    \node[left=11pt] at (q.center) {$Q$};
    \node[right=2pt] at (q.center) {$R_{\mathrm{on}}$};
    \node[below=5pt] at (diode.center) {$R_{\mathrm{D}},\ V_{\mathrm{D}}$};
    \node[below=5pt] at (1.65,2.7) {$R_{\mathrm{L}}$};
    \draw (6.0,2.7) to[C,l_=$C$,i>^=$i_{\mathrm{C}}$] (6.0,0);
    \draw (7.5,2.7) to[european resistor,l_=$R$] (7.5,0) -- (0,0);
    \node at (8.05,2.7) {$+$};
    \node at (8.05,0) {$-$};
    \node at (8.05,1.35) {$v_{\mathrm{out}}$};
  \end{notescircuit}
  \vspace{-0.15cm}
  \caption{Boost 电路的电感与半导体通态损耗参数}
  \label{pic0209}
\end{figure}
}

\DefineFigure{pic0210}{
\begin{figure}[H]
  \vspace{-0.2cm}
  \centering
  \begin{subfigure}[t]{0.48\textwidth}
    \centering
    \begin{notescircuit}[scale=0.9]
      \ctikzset{
        tripoles/nigfete/base width=0.32,
        tripoles/nigfete/gate width=0.40
      }
      \coordinate (gnd-boost-loss-original) at (0,0);
      \node[nigfete, nogate, anchor=center] (q-boost-loss-original) at (2.94,0.95) {};
      \draw (gnd-boost-loss-original)
      to[V, l_=$V_{\mathrm{in}}$, name=vin-boost-loss-original] (0,2.0)
      to[L, l=$L$, name=ind-boost-loss-original] (2.94,2.0) coordinate (sw-boost-loss-original)
      to[D-, l=$D$, name=diode-boost-loss-original] (4.84,2.0) coordinate (out-boost-loss-original)
      -- ++(1.05,0) coordinate (load-boost-loss-original);
      \draw[-{Latex[length=1.5mm]}] (0.12,2.0) -- (0.52,2.0) node[midway,above] {$i_{\mathrm{L}}$};
      \draw (sw-boost-loss-original) -- (q-boost-loss-original.D);
      \draw (q-boost-loss-original.S) -- (q-boost-loss-original.S |- gnd-boost-loss-original);
      \draw (q-boost-loss-original.inner down -| q-boost-loss-original.nogate) -- ++(-0.22,0);
      \node[left=11pt] at (q-boost-loss-original.center) {$Q$};
      \node[right=2pt] at (q-boost-loss-original.center) {$R_{\mathrm{on}}$};
      \node[below=5pt] at (diode-boost-loss-original.center) {$R_{\mathrm{D}},\ V_{\mathrm{D}}$};
      \node[left=1pt] at (0,1.35) {$+$};
      \node[left=1pt] at (0,0.65) {$-$};
      \node[below=5pt] at (ind-boost-loss-original.center) {$R_{\mathrm{L}}$};
      \draw (out-boost-loss-original) to[C, l_=$C$] (out-boost-loss-original |- gnd-boost-loss-original);
      \draw[-{Latex[length=1.5mm]}] ($(out-boost-loss-original)+(0,-0.08)$) -- ($(out-boost-loss-original)+(0,-0.52)$)
      node[midway,right] {$i_{\mathrm{C}}$};
      \draw (load-boost-loss-original) to[european resistor, l_=$R$] (load-boost-loss-original |- gnd-boost-loss-original);
      \draw (load-boost-loss-original |- gnd-boost-loss-original) -- (gnd-boost-loss-original);
      \node[right] at (load-boost-loss-original) {$+$};
      \node[right] at (load-boost-loss-original |- gnd-boost-loss-original) {$-$};
      \path (load-boost-loss-original) -- node[right] {$v_{\mathrm{out}}$} (load-boost-loss-original |- gnd-boost-loss-original);
    \end{notescircuit}\par
    \caption{原电路}
    \label{pic0210a}
  \end{subfigure}
  \hfill
  \begin{subfigure}[t]{0.48\textwidth}
    \centering
    \begin{notescircuit}[scale=0.9]
      \coordinate (gnd-boost-loss-switch) at (0,0);
      \coordinate (pole-boost-loss-switch) at (2.94,2.0);
      \coordinate (throw-one-boost-loss-switch) at (3.24,1.25);
      \coordinate (throw-two-boost-loss-switch) at (3.54,2.0);
      \draw (gnd-boost-loss-switch)
      to[V, l_=$V_{\mathrm{in}}$, name=vin-boost-loss-switch] (0,2.0)
      -- (0.56,2.0)
      to[L, l=$L$, name=ind-boost-loss-switch] (1.64,2.0)
      to[european resistor, l=$R_{\mathrm{L}}$] (pole-boost-loss-switch);
      \draw[-{Latex[length=1.5mm]}] (0.12,2.0) -- (0.52,2.0) node[midway,above] {$i_{\mathrm{L}}$};
      \node[ocirc] at (pole-boost-loss-switch) {};
      \node[ocirc] at (throw-one-boost-loss-switch) {};
      \node[ocirc] at (throw-two-boost-loss-switch) {};
      \draw (pole-boost-loss-switch) -- (3.44,1.60);
      \node[left] at (throw-one-boost-loss-switch) {$1$};
      \node[above] at (throw-two-boost-loss-switch) {$2$};
      \draw (throw-one-boost-loss-switch)
      to[european resistor, l_=$R_{\mathrm{on}}$] (throw-one-boost-loss-switch |- gnd-boost-loss-switch)
      -- (gnd-boost-loss-switch);
      \draw (throw-two-boost-loss-switch)
      -- (3.84,2.0)
      to[european voltage source, l=$V_{\mathrm{D}}$, name=diode-drop-boost-loss-switch] (4.74,2.0)
      to[european resistor, l=$R_{\mathrm{D}}$] (5.54,2.0)
      -- (5.84,2.0) coordinate (out-boost-loss-switch)
      -- (6.59,2.0) coordinate (load-boost-loss-switch);
      \node[left=1pt] at (0,1.35) {$+$};
      \node[left=1pt] at (0,0.65) {$-$};
      \node[below=4pt] at (ind-boost-loss-switch.center)
      {$+\;v_{\mathrm{L}}\;-$};
      \node[anchor=base] at ($(diode-drop-boost-loss-switch.center)+(-0.35,-0.35)$) {$+$};
      \node[anchor=base] at ($(diode-drop-boost-loss-switch.center)+(0.35,-0.35)$) {$-$};
      \draw (out-boost-loss-switch) to[C, l_=$C$] (out-boost-loss-switch |- gnd-boost-loss-switch);
      \draw[-{Latex[length=1.5mm]}] ($(out-boost-loss-switch)+(0,-0.08)$) -- ($(out-boost-loss-switch)+(0,-0.52)$)
      node[midway,right] {$i_{\mathrm{C}}$};
      \draw (load-boost-loss-switch) to[european resistor, l_=$R$] (load-boost-loss-switch |- gnd-boost-loss-switch);
      \draw (load-boost-loss-switch |- gnd-boost-loss-switch) -- (gnd-boost-loss-switch);
      \node[right] at (load-boost-loss-switch) {$+$};
      \node[right] at (load-boost-loss-switch |- gnd-boost-loss-switch) {$-$};
      \path (load-boost-loss-switch) -- node[right] {$v_{\mathrm{out}}$} (load-boost-loss-switch |- gnd-boost-loss-switch);
    \end{notescircuit}\par
    \caption{等效电路}
    \label{pic0210b}
  \end{subfigure}
  \medskip
  \begin{subfigure}[t]{0.48\textwidth}
    \centering
    \begin{notescircuit}[scale=0.9]
      \coordinate (gnd-boost-loss-a) at (0,0);
      \draw (gnd-boost-loss-a)
      to[V, l_=$V_{\mathrm{in}}$, name=vin-boost-loss-a] (0,1.8)
      -- (0.56,1.8)
      to[L, l=$L$, name=ind-boost-loss-a] (1.64,1.8)
      to[european resistor, l=$R_{\mathrm{L}}$] (2.94,1.8) coordinate (sw-boost-loss-a)
      to[european resistor, l_=$R_{\mathrm{on}}$] (sw-boost-loss-a |- gnd-boost-loss-a)
      -- (gnd-boost-loss-a);
      \draw[-{Latex[length=1.5mm]}] (0.12,1.8) -- (0.52,1.8) node[midway,above] {$i_{\mathrm{L}}$};
      \node[left=1pt] at (0,1.25) {$+$};
      \node[left=1pt] at (0,0.55) {$-$};
      \node[below=4pt] at (ind-boost-loss-a.center)
      {$+\;v_{\mathrm{L}}\;-$};
      \draw (3.84,1.8) coordinate (out-boost-loss-a)
      -- ++(1.05,0) coordinate (load-boost-loss-a);
      \draw (out-boost-loss-a) to[C, l_=$C$] (out-boost-loss-a |- gnd-boost-loss-a);
      \draw[-{Latex[length=1.5mm]}] ($(out-boost-loss-a)+(0,-0.08)$) -- ($(out-boost-loss-a)+(0,-0.52)$)
      node[midway,right] {$i_{\mathrm{C}}$};
      \draw (load-boost-loss-a) to[european resistor, l_=$R$] (load-boost-loss-a |- gnd-boost-loss-a);
      \draw (load-boost-loss-a |- gnd-boost-loss-a) -- (gnd-boost-loss-a);
      \node[right] at (load-boost-loss-a) {$+$};
      \node[right] at (load-boost-loss-a |- gnd-boost-loss-a) {$-$};
      \path (load-boost-loss-a) -- node[right] {$v_{\mathrm{out}}$} (load-boost-loss-a |- gnd-boost-loss-a);
    \end{notescircuit}\par
    \caption{状态 1}
    \label{pic0210c}
  \end{subfigure}
  \hfill
  \begin{subfigure}[t]{0.48\textwidth}
    \centering
    \begin{notescircuit}[scale=0.9]
      \coordinate (gnd-boost-loss-b) at (0,0);
      \draw (gnd-boost-loss-b)
      to[V, l_=$V_{\mathrm{in}}$, name=vin-boost-loss-b] (0,1.8)
      -- (0.56,1.8)
      to[L, l=$L$, name=ind-boost-loss-b] (1.64,1.8)
      to[european resistor, l=$R_{\mathrm{L}}$] (2.94,1.8)
      to[european voltage source, l=$V_{\mathrm{D}}$, name=diode-drop-boost-loss-b] (3.94,1.8)
      to[european resistor, l=$R_{\mathrm{D}}$] (5.24,1.8) coordinate (out-boost-loss-b)
      -- ++(1.05,0) coordinate (load-boost-loss-b);
      \draw[-{Latex[length=1.5mm]}] (0.12,1.8) -- (0.52,1.8) node[midway,above] {$i_{\mathrm{L}}$};
      \node[left=1pt] at (0,1.25) {$+$};
      \node[left=1pt] at (0,0.55) {$-$};
      \node[below=4pt] at (ind-boost-loss-b.center)
      {$+\;v_{\mathrm{L}}\;-$};
      \node[anchor=base] at ($(diode-drop-boost-loss-b.center)+(-0.35,-0.35)$) {$+$};
      \node[anchor=base] at ($(diode-drop-boost-loss-b.center)+(0.35,-0.35)$) {$-$};
      \draw (out-boost-loss-b) to[C, l_=$C$] (out-boost-loss-b |- gnd-boost-loss-b);
      \draw[-{Latex[length=1.5mm]}] ($(out-boost-loss-b)+(0,-0.08)$) -- ($(out-boost-loss-b)+(0,-0.52)$)
      node[midway,right] {$i_{\mathrm{C}}$};
      \draw (load-boost-loss-b) to[european resistor, l_=$R$] (load-boost-loss-b |- gnd-boost-loss-b);
      \draw (load-boost-loss-b |- gnd-boost-loss-b) -- (gnd-boost-loss-b);
      \node[right] at (load-boost-loss-b) {$+$};
      \node[right] at (load-boost-loss-b |- gnd-boost-loss-b) {$-$};
      \path (load-boost-loss-b) -- node[right] {$v_{\mathrm{out}}$} (load-boost-loss-b |- gnd-boost-loss-b);
    \end{notescircuit}\par
    \caption{状态 2}
    \label{pic0210d}
  \end{subfigure}
  \vspace{-0.5cm}
  \caption{含电感铜损及半导体通态损耗的 Boost 变换器及其开关状态等效电路}
  \vspace{-0.2cm}
  \label{pic0210}
\end{figure}
}

\DefineFigure{pic0211}{
\begin{figure}[H]
  \vspace{-0.3cm}
  \centering
  \begin{subfigure}[t]{0.48\textwidth}
    \centering
    \begin{notescircuit}[scale=0.7]
      \ctikzset{sources/symbol/sign rotation=straight}
      \draw (0,2.6) to[european voltage source,l=$V_{\mathrm{in}}$,name=loss-source-model] (0,0);
      \node[above left=1pt] at (loss-source-model.b) {$+$};
      \node[below left=1pt] at (loss-source-model.a) {$-$};
      \draw (0,2.6) -- (0.9,2.6)
      to[european resistor,l=$R_{\Sigma}$] (2.2,2.6)
      to[european voltage source,l=$D'V_{\mathrm{D}}$,name=loss-drop-model] (3.2,2.6)
      -- (4.4,2.6);
      \node[anchor=base] at ($(loss-drop-model.center)+(-0.5,-0.45)$) {$+$};
      \node[anchor=base] at ($(loss-drop-model.center)+(0.5,-0.45)$) {$-$};
      \draw (4.4,2.6) to[american controlled voltage source,l_=$D'V_{\mathrm{out}}$] (4.4,0);
      \draw (4.4,0) -- (0,0);
      \draw[-{Latex[length=1.7mm]}] (3.35,2.6) -- (4.05,2.6) node[midway,above] {$I_{\mathrm{L}}$};
      \draw (5.8,0) to[american controlled current source,l_=$D'I_{\mathrm{L}}$] (5.8,2.6)
      -- (8.1,2.6) to[european resistor,l_=$R$] (8.1,0) -- (5.8,0);
      \draw[-{Latex[length=1.7mm]}] (6.3,2.6) -- (7.2,2.6) node[midway,above] {$I_{\mathrm{out}}$};
      \node at (8.65,2.6) {$+$};
      \node at (8.65,0) {$-$};
      \node at (8.65,1.3) {$V_{\mathrm{out}}$};
    \end{notescircuit}\par
    \caption{受控源模型}
    \label{pic0211a}
  \end{subfigure}
  \hfill
  \begin{subfigure}[t]{0.48\textwidth}
    \centering
    \begin{notescircuit}[scale=0.74]
      \node[transformer] (tr) at (4.0,1.3) {};
      \node[circ] at (tr.outer dot A1) {};
      \node[circ] at (tr.outer dot B1) {};
      \node[above] at (4.0,3.05) {$D':1$};
      \coordinate (loss-dc-mark) at ($(tr-L1.a)!{0.5/\ctikzvalof{bipoles/americaninductor/coils}}!(tr-L1.b)$);
      \draw (tr.outer dot A1 |- loss-dc-mark) -- (tr.outer dot B1 |- loss-dc-mark);
      \draw (0,2.6) to[european voltage source,l=$V_{\mathrm{in}}$,name=loss-source-transformer] (0,0);
      \node[above left=1pt] at (loss-source-transformer.b) {$+$};
      \node[below left=1pt] at (loss-source-transformer.a) {$-$};
      \draw (0,2.6) to[european resistor,l=$R_{\Sigma}$] (1.3,2.6)
      to[european voltage source,l=$D'V_{\mathrm{D}}$,name=loss-drop-transformer] (2.3,2.6)
      -- (3.0,2.6) -| (tr.A1);
      \node[anchor=base] at ($(loss-drop-transformer.center)+(-0.5,-0.45)$) {$+$};
      \node[anchor=base] at ($(loss-drop-transformer.center)+(0.5,-0.45)$) {$-$};
      \draw (tr.A2) |- (0,0);
      \draw (tr.B1) |- (6.3,2.6);
      \draw (6.3,2.6) to[european resistor,l_=$R$] (6.3,0) -| (tr.B2);
      \node at (6.85,2.6) {$+$};
      \node at (6.85,0) {$-$};
      \node at (6.85,1.3) {$V_{\mathrm{out}}$};
      \draw[-{Latex[length=1.7mm]}] (2.4,2.6) -- (3.1,2.6) node[midway,above] {$I_{\mathrm{L}}$};
      \draw[-{Latex[length=1.7mm]}] (5.0,2.6) -- (5.7,2.6) node[midway,above] {$I_{\mathrm{out}}$};
      \draw[-{Latex[length=1.7mm]}] (4.0,-0.35) -- ([yshift=-0.05cm]tr.south) node[midway,left] {$D$};
    \end{notescircuit}\par
    \caption{直流“变压器”模型}
    \label{pic0211b}
  \end{subfigure}
  \vspace{-0.15cm}
  \caption{含电感铜损及半导体通态损耗的 Boost 变换器直流稳态等效模型}
  \vspace{-0.3cm}
  \label{pic0211}
\end{figure}
}

\DefineFigure{pic0212}{
\begin{figure}[H]
  \vspace{-0.3cm}
  \centering
  \begin{notescircuit}[scale=0.95]
    \draw (0,2.6) to[european voltage source,l=$V_{\mathrm{in}}/D'-V_{\mathrm{D}}$,name=loss-source-reflected] (0,0);
    \node[above left=1pt] at (loss-source-reflected.b) {$+$};
    \node[below left=1pt] at (loss-source-reflected.a) {$-$};
    \draw (0,2.6) -- (0.4,2.6)
    to[european resistor,l=$R_{\Sigma}/(D')^2$] (1.7,2.6) -- (4.0,2.6)
    to[european resistor,l_=$R$] (4.0,0) -- (0,0);
    \draw[-{Latex[length=1.7mm]}] (2.1,2.6) -- (3.2,2.6) node[midway,above] {$D'I_{\mathrm{L}}$};
    \node at (4.55,2.6) {$+$};
    \node at (4.55,0) {$-$};
    \node at (4.55,1.3) {$V_{\mathrm{out}}$};
  \end{notescircuit}
  \vspace{-0.15cm}
  \caption{含电感铜损及半导体通态损耗的 Boost 稳态模型折算到输出侧}
  \vspace{-0.3cm}
  \label{pic0212}
\end{figure}
}
